\chapter{The modified force field}
\label{ch:ff}

A force field describes a molecule held together by chemical bonds.
Once its atoms are ionized, the electrons that formed those bonds are gone, and the bonded terms would hold the
exploding fragments together with a strength they no longer have.
With \code{mdrun -ionize} the bonded and 1--4 interactions are therefore modified as the atoms charge up
(\file{src/gmxlib/bondfree.c} and \file{src/gmxlib/nonbonded/nonbonded.c}; the switch is the global
\code{moldstruct\_altered\_ff}).
Without \code{-ionize} the force field is that of stock GROMACS.

\section{Charges}

At the first step the module replaces every atom's partial charge by its ionization charge $q = Z-(K+L+M)$, zero for
a neutral atom.
From then on the Coulomb interactions act only between ions, with integer charges that change as the atoms ionize.
The Lennard-Jones interactions are not modified.

\section{The ionization factor}

Every modified bonded term is scaled by the \emph{ionization factor} of the mean charge $\bar q$ of the atoms it
involves:
\begin{equation}
  f(\bar q) = \min\!\big(1,\ \max(1-\bar q,\ 0)\big),
\end{equation}
which is 1 for neutral atoms, falls linearly, and is 0 once the atoms carry on average one charge.
Energies and forces are scaled together, so the scaled terms remain conservative at fixed charge.

\begin{table}[H]
\centering\small
\begin{tabularx}{\textwidth}{@{}l l X@{}}
\toprule
interaction & GROMACS function & modification under \code{-ionize} \\
\midrule
Morse bond & \code{morse\_bonds} & breaks irreversibly at $\beta(r-b_0) > 4$; then removed and replaced by the pair's
  $C_{12}/r^{12}$ repulsion; the bare 1--2 Coulomb between the two ions is added \\
harmonic bond & \code{bonds} & not modified --- use \code{morse = yes} \\
angle & \code{angles} & scaled by $f(\bar q)$ of the three atoms; explicit 1--3 Coulomb added \\
Urey--Bradley & \code{urey\_bradley} & angle part scaled by $f(\bar q)$; the 1--3 spring breaks irreversibly at
  $r_{13} > 2\,r_{13}^0$, scaled by $f(\bar q)$ before; explicit 1--3 Coulomb added \\
proper dihedral & \code{pdihs} & scaled by $f(\bar q)$ of the four atoms \\
improper dihedral & \code{idihs} & scaled by $f(\bar q)$ of the four atoms \\
Ryckaert--Bellemans & \code{rbdihs} & scaled by $f(\bar q)$ of the four atoms \\
CMAP & \code{cmap\_dihs} & scaled by a stepped $f(\bar q)$ of the five atoms: 1, 0.8, 0.6, 0.4, 0.2, 0 at
  $\bar q <$ 0.19, 0.39, 0.59, 0.79, 1, $\ge1$ \\
1--4 pair & \code{do\_listed\_vdw\_q} & Coulomb scaled from \code{fudgeQQ} to full strength, LJ unchanged \\
\bottomrule
\end{tabularx}
\caption{The bonded and 1--4 interactions under \code{-ionize}.}
\label{tab:ff}
\end{table}

\section{Bond breaking}

\paragraph{Morse bonds.}
The bonds should be Morse potentials, $V = D\,[1-e^{-\beta(r-b_0)}]^2$, which \code{grompp} makes from harmonic bonds with
\code{morse = yes} (dissociation energies from \file{share/top/edissoc.dat}); the example does this.
A Morse bond breaks permanently once $\beta(r-b_0) > 4$.
The threshold is in units of the decay length $1/\beta$, so it lies at the same point of the dissociation curve for
every bond type: about 7\,\% of the peak restoring force is left and 96\,\% of the well depth has been paid.
(A multiple of $b_0$ would be much more aggressive for short bonds: $2b_0$ leaves 43\,\% of the peak force on a C--H bond,
20\,\% on C--C.)
After breaking, the Morse term is gone.
Because GROMACS's exclusion lists are fixed, the pair stays excluded from the non-bonded interactions; the bond term
therefore puts back the pair's own repulsive wall $C_{12}/r^{12}$ (the real $C_{12}$ of the two atom types, without the
$r^{-6}$ dispersion), so that broken fragments cannot pass through each other.

\paragraph{The 1--2 Coulomb interaction.}
Bonded pairs are excluded from the non-bonded Coulomb sum, which is right for partial charges in a molecule but wrong for
two ions.
The Morse term therefore adds the bare Coulomb interaction of the pair, $\kc q_iq_j/r$, broken or not; its energy is
reported in the Morse term of the energy file.

\paragraph{Urey--Bradley 1--3 springs} (CHARMM) are harmonic and have no decay length, so their threshold is
geometric: the spring breaks permanently at twice its equilibrium length, where the angle it braces is already
destroyed.

\paragraph{Harmonic bonds} are not modified: they neither weaken nor break.
Runs with \code{-ionize} should use \code{morse = yes}.

\paragraph{The broken-bond registry.}
Broken pairs are stored by their global atom indices in a hash set, separately for Morse and Urey--Bradley pairs.
It is thread-local: under particle decomposition each rank owns a fixed slice of the bonded interactions, so a bond
is only ever evaluated on one rank.
It is not written to the checkpoint; after a restart a bond still beyond its threshold re-registers at the first
step, so only bonds that broke and came back inside the threshold are forgotten.

\section{Angles and the 1--3 interaction}

Angles are scaled by the ionization factor of their three atoms.
As with bonds, 1--3 pairs are excluded from the non-bonded Coulomb sum, so the angle (and Urey--Bradley) terms add
the 1--3 Coulomb interaction explicitly, as a genuine potential: energy and force, with the signed charge product.

\section{The 1--4 interactions}

The 1--4 Coulomb interaction of a force field is scaled by \code{fudgeQQ} (e.g.\ 0.5--1, depending on the force
field) because the bonded terms already describe part of it.
Once the bonded terms fade, the 1--4 Coulomb is interpolated back to full strength with the mean charge of the pair:
\begin{equation}
  \epsilon_{14} = \epsilon\,\big[1 - (1-\code{fudgeQQ})\,f(\bar q_{ij})\big],
\end{equation}
\code{fudgeQQ} for neutral atoms and 1 for ionized ones.
The 1--4 Lennard-Jones term is not changed.
